% No preamble. Prefix ccs-. Actual covariant moving-reference tangents.
\section{A covariant moving-reference scalar average}
\label{ccs-section}

We construct tangent operators for the actual ordinary composite
average. The nonlinear constraint and its cutoff are unchanged.
The construction uses the internal group-valued transport selected
in that ordinary map, rather than an unverified identification with
an older axial operator carrying the name \(Q'\).

\subsection{Actual data and the primitive relative correction}
\label{ccs-primitive-section}

Fix odd \(\ell\ge3\), \(K=\ell^k\), and compatible nested centered
blocks. Coordinates are current-link angles with counting scalar
and bond inner products. Use
\[
 U^g(e)=g(e_-)U(e)g(e_+)^{-1},\qquad
 (D_U\phi)(e)=\phi(e_-)-\operatorname{Ad}_{U(e)}\phi(e_+).
\]
The generator acts on all scalar sites. Scalar parameters have
angle units at every level; no spacing factor is inserted into
this formula.

Let \(P_\ell(U;y,x)\) be exactly the internal group-valued transport
entering the selected ordinary map \(M_\ell\). The actual definition
in Ba{\l}aban, CMP109 (1987), pp.~252--254, (0.3),(0.5)--(0.12),
supplies local analyticity, internal-block support, and
\begin{equation}
 P_\ell(U^g;y,x)=g(y)P_\ell(U;y,x)g(x)^{-1}.
 \label{ccs-endpoint}
\end{equation}
At identity its derivative is the same arithmetic contour mean
that occurs in the ordinary arithmetic differential.
These facts also hold when \(P_\ell\) is the selected group mean
of coordinate-order transporters; it need not be one axial path.

The pinned deposition specifies this ordinary map at Raw I:
831--845. Its relative contour and ordinary increments occur at
46965--46978, and the moving-reference endpoint correction at
47076--47086. Raw I:23913--23920 describes an older axial/nested
scalar average. No equality with every such convention is needed:
the following operators are defined directly using the actual
\(P_\ell\).

Freeze \(U\), put \(W=M_\ell(U)\), and for a small bond increment
\(x\) define
\begin{equation}
 \begin{split}
 r_U(x;y,z)
   &=\frac1i\log\!\left[
       P_\ell(e^{ix}U;y,z)P_\ell(U;y,z)^{-1}\right],\\
 \alpha_U(x;y)&=\ell^{-d}\sum_{z\in B_\ell(y)}r_U(x;y,z),\\
 a_U(x;y)&=e^{i\alpha_U(x;y)}.
 \end{split}
 \label{ccs-relative-mean}
\end{equation}
All relative logarithms vanish at \(x=0\), and \(a_U(0)=1\).
Let
\begin{equation}
 Z_Ux=D\alpha_U(0)[x],\qquad
 C_Ux=\frac1i\left.\frac d{ds}\right|_0
                  M_\ell(e^{isx}U)W^{-1}.
 \label{ccs-tangents}
\end{equation}
The local corrected output with this frozen reference is
\[
 \overline F_U(x;c)=a_U(x;c_-)^{-1}
                         M_\ell(e^{ix}U)(c)a_U(x;c_+).
\]
It equals \(W\) at zero. Its left-relative differential is exactly
\begin{equation}
                     \overline C_U=C_U-D_WZ_U.
 \label{ccs-primitive-correction}
\end{equation}
This is a moving-reference tangent identity. It is not an
identification with the derivative of the fixed-identity corrected
map at \(U\).

\begin{lemma}[Primitive scalar intertwining]
\label{ccs-primitive-lemma}
Define
\begin{equation}
 (B_U\phi)(y)=\ell^{-d}\sum_{z\in B_\ell(y)}
                  \operatorname{Ad}_{P_\ell(U;y,z)}\phi(z).
 \label{ccs-primitive-B}
\end{equation}
With \(r_\ell\) restriction to the new centers,
\begin{equation}
 Z_UD_U=r_\ell-B_U,\qquad
                  \overline C_UD_U=D_WB_U.
 \label{ccs-primitive-intertwining}
\end{equation}
\end{lemma}

\begin{proof}
For the fine gauge \(e^{is\phi}\), the relative contour derivative
from \eqref{ccs-endpoint} is
\[
 \phi(y)-\operatorname{Ad}_{P_\ell(U;y,z)}\phi(z).
\]
Its fine left tangent is \(D_U\phi\). Averaging this expression
proves the first identity. There is no logarithm Jacobian at this
step: its argument is identity at the reference point.
Ordinary covariance gives \(C_UD_U=D_Wr_\ell\). Subtract
\(D_WZ_UD_U\) using \eqref{ccs-primitive-correction}.
\end{proof}

\subsection{Joint covariance and composition}
\label{ccs-composition-section}

For a fixed site gauge \(g\), let \(O_E(g),O_S(g),O_c(g)\)
denote adjoint frames on fine bond source sites, fine scalar sites,
and new centers. Since
\[
 e^{iO_E(g)x}U^g=(e^{ix}U)^g,
\]
the relative contour product at \(U^g\) is conjugated by \(g(y)\).
The similarity-equivariant logarithm therefore gives
\begin{equation}
 \begin{split}
 Z_{U^g}O_E(g)&=O_c(g)Z_U,\\
 \overline C_{U^g}O_E(g)&=O_{E_c}(g)\overline C_U,\\
 B_{U^g}O_S(g)&=O_c(g)B_U.
 \end{split}
 \label{ccs-primitive-covariance}
\end{equation}
Both the reference and the varied configuration are transformed.
No multiplicative cocycle in the reference is assumed.
Adjoint frames are orthogonal for the real invariant trace pairing,
and complex orthogonal for its holomorphic bilinear continuation.

At stage \(j\), evaluate \(C_j,\overline C_j,Z_j,B_j,D_j\)
at the ordinary reference
\[
                         U_j=M_\ell^j(U).
\]
Set
\begin{equation}
 \begin{split}
 C^{(t)}&=C_{t-1}\cdots C_0,\\
 \overline C^{(t)}&=\overline C_{t-1}\cdots\overline C_0,\\
 B^{(t)}&=B_{t-1}\cdots B_0,
 \end{split}
 \label{ccs-products}
\end{equation}
with identity products at \(t=0\).
The first is exactly the derivative of the actual ordinary
composition. Define \(Z^{(0)}=0\) and
\begin{equation}
        Z^{(t+1)}=B_tZ^{(t)}+Z_tC^{(t)}.
 \label{ccs-Z-recursion}
\end{equation}

\begin{proposition}[Exact composite tangent identities]
\label{ccs-composite-theorem}
At the final level, writing \(C,\overline C,B,Z,D_c\) for these
products and the final generator,
\begin{equation}
 C-\overline C=D_cZ,\qquad
 \overline C D_{\rm all}=D_cB,\qquad
                  ZD_{\rm all}=r_K-B.
 \label{ccs-composite-identities}
\end{equation}
The operators have exact joint gauge covariance with the original
fine and final ordinary reference frames.
\end{proposition}

\begin{proof}
The primitive intertwining gives
\(\overline C^{(t)}D_0=D_tB^{(t)}\).
Suppose \(C^{(t)}=\overline C^{(t)}+D_tZ^{(t)}\).
Then
\[
 C_tC^{(t)}
  =\overline C_t\overline C^{(t)}
        +D_{t+1}\{B_tZ^{(t)}+Z_tC^{(t)}\},
\]
which proves the first two identities by induction.
For the third, insert
\(Z_tD_t=r_\ell-B_t\) and
\(C^{(t)}D_0=D_tr_{\ell^t}\) in
\eqref{ccs-Z-recursion}; the two intermediate
\(B_tr_{\ell^t}\) terms cancel.
Finally, \eqref{ccs-primitive-covariance} telescopes at the
ordinary reference backgrounds. These are not the
fixed-identity corrected iterates.
\end{proof}

\subsection{Own-block locality and an exact constant row Gram matrix}
\label{ccs-locality-section}

For \(x\in B_K(y)\), let
\(x=x_0,x_1,\ldots,x_k=y\) be its unique nested sequence of
block centers. The scalar product in \eqref{ccs-products}
has only this one ancestry. Consequently
\begin{equation}
 \begin{split}
 B(U;y,x)&=K^{-d}\operatorname{Ad}_{P_{\rm hier}(U;y,x)},\\
 P_{\rm hier}(U;y,x)
  &=P_\ell(U_{k-1};x_k,x_{k-1})\cdots
                                 P_\ell(U_0;x_1,x_0),
 \end{split}
 \label{ccs-hierarchical-B}
\end{equation}
and \(B(U;y,x)=0\) outside \(B_K(y)\).
Each factor is one group element, even when obtained from a
selected group mean of path transports.

Its background dependencies lie in the same final block.
A stage contour uses internal stage bonds. Each such ordinary
bond reads the two descendant blocks at its endpoints, both
inside the parent block by centered nested tiling. Iterating
downwards proves that no coefficient of this scalar row reads
a fine background edge crossing its final block boundary.
The same conclusion holds for all of its derivatives.

The rows of \(C\) and \(\overline C\) read only the two final
endpoint blocks. The scalar \(Z\) row stays in one block:
in \eqref{ccs-Z-recursion}, \(Z_t\) reads internal stage bonds,
and \(C^{(t)}\) on these bonds reads only descendants inside
their parent block. This is full coefficient locality, not
merely support at flatness.

\begin{theorem}[Covariant block projection]
\label{ccs-projector-theorem}
For every real or complex background in the common germ,
\begin{equation}
                 BB^T=K^{-d}I.
 \label{ccs-row-Gram}
\end{equation}
In dimension four,
\begin{equation}
                  P_B=I-K^4B^TB
 \label{ccs-projector}
\end{equation}
is the symmetric projector onto \(\ker B\), orthogonal for
real fields. It is holomorphic, has one-final-block matrix
and background support, and obeys
\begin{equation}
          P_B(U^g)=O_S(g)P_B(U)O_S(g)^T.
 \label{ccs-projector-covariance}
\end{equation}
The normalized term \(\frac12\log\det BB^T\) vanishes identically.
\end{theorem}

\begin{proof}
Distinct rows have disjoint scalar supports because the final
blocks partition the fine sites. In the invariant bilinear
pairing,
\(\operatorname{Ad}_P\operatorname{Ad}_P^T=I\).
Equation \eqref{ccs-hierarchical-B} therefore gives
\[
 (BB^T)_{yy}
      =K^{-2d}\sum_{x\in B_K(y)}I=K^{-d}I,
\]
with zero off-diagonal blocks. This algebra remains valid on
the complexified group. The projector assertions follow
directly, and covariance follows from that of \(B\).
\end{proof}

\subsection{Flat values and the analytic domain}
\label{ccs-analytic-section}

At identity, the contour derivative in \(Z\) is the root term
of the actual ordinary arithmetic differential. Subtraction
of the two endpoint terms in
\eqref{ccs-primitive-correction} leaves precisely
\begin{equation}
 (\overline C_1x)(c)=\ell^{-d}
       \sum_{z\in B_\ell(c_-)}x([z,z+\ell e_\mu]),
 \qquad B_1=Q_{0,\ell}.
 \label{ccs-flat-primitive}
\end{equation}
Centered-block tiling then gives
\begin{equation}
             \overline C(1)=A_K^{\rm st},\qquad
                         B(1)=Q_{0,K}.
 \label{ccs-flat-composite}
\end{equation}
The bond map here is an angle map, with a primitive factor
\(\ell\) relative to the normalized physical-potential map.
There is no such factor for the scalar average.

At fixed \(K\), each coefficient is a composition of local
analytic maps reading finitely many links in one or two final
blocks. A small link-angle product ball therefore works uniformly
in volume and in unrelated preceding depth \(M\).
Its radius and constants may depend on \(K\).
Cauchy derivatives preserve the same support. Compatible tori
are read with the prescribed local lifts and nested partition.

The previously proved ordinary-reference trajectory gives a
stronger optional bound. For \(U=e^{iA/K}\) in its small globally
bounded physical-potential ball, stage angles are
\(O(\ell^j\|A\|_\infty/K)\). Fixed-\(\ell\) local analytic
Schur bounds and the sum
\(\sum_{j<k}\ell^j/K\le(\ell-1)^{-1}\) imply
\begin{equation}
 \|B-Q_{0,K}\|_{{\rm count},\,2\to2}
                    \le C\|A\|_\infty K^{-2}
                    \qquad(d=4).
 \label{ccs-physical-smallness}
\end{equation}
The corresponding counting scale for the angle bond map
\(\overline C\) is \(K^{-1}\).
These stronger estimates are not needed for fixed-\(K\)
volume uniformity or the exact identity
\eqref{ccs-row-Gram}.

The induced scalar map of the fixed-identity corrected
composition is a different operator. It contains logarithm
Jacobians, and its row Gram matrix need not be constant.
The present moving-reference construction avoids that issue;
it must not be replaced by the fixed-identity derivative
without the appropriate endpoint conversion.

\subsection{A basis-free scalar completion}
\label{ccs-completion-section}

Let \(G_{\rm all}=D_{\rm all}^TD_{\rm all}\), and let \(V\)
be any bilinear orthonormal frame of \(\ker B\). Define
\begin{equation}
 G_B=V^TG_{\rm all}V,\qquad
 \widehat G_B=P_BG_{\rm all}P_B+K^{-2}(I-P_B).
 \label{ccs-completion}
\end{equation}
The two projector subspaces reduce \(\widehat G_B\); hence
\begin{equation}
       \det\widehat G_B=\det G_B\,
                        (K^{-2})^{\operatorname{rank}B}.
 \label{ccs-completion-determinant}
\end{equation}
The extra factor is background independent.
The ambient matrix \(\widehat G_B\) is exactly gauge covariant,
without a need to choose gauge-covariant frames \(V\).

For the companion vector matrix
\[
 Q_B=\mathcal H+D_{\rm all}P_BD_{\rm all}^T
                         +a\overline C^T\overline C,
\]
the actual residual covariance and
\eqref{ccs-projector-covariance} give the same conclusion.
When the general tangent determinant identity and the needed
spectral gaps are supplied separately, these observations
permit a formulation using the two ambient matrices
\(\widehat G_B,Q_B\), with no isolated global base site and
no normalized row-Gram term.

Locality and covariance alone do not prove those spectral gaps
or the determinant identity. Once they hold, a fixed-\(K\)
uniform complex perturbation neighborhood permits the usual
diagonal logarithm argument. This does not compute a
large-\(K\) vector Maxwell coefficient or establish a
continuum Yang--Mills theory.
